\chapter{Requisitos}\label{cap:requisitos}

Este capítulo reúne lo que debe estar disponible antes de construir Qadra
desde sus fuentes y antes de ejecutarlo en un servidor. Las versiones
indicadas son las que fija el repositorio; cuando no se cita otro archivo,
los requisitos de ejecución proceden de \archivo{docs/deployment.md}. El
único entorno soportado para construir y ejecutar es
Linux x86\_64 (\archivo{docs/document-format-operations.md}).

\section{Construcción y ejecución}\label{sec:req-tabla}

La \cref{tab:req-software} separa dos situaciones. La \emph{construcción}
ocurre en el equipo de desarrollo o en el ejecutor de integración continua:
compila el binario \texttt{despacho-cli}, los módulos Rust/WASM y los dos
portales web. La \emph{ejecución} ocurre en el servidor, que recibe un
paquete ya compilado y no necesita las herramientas de construcción. La
construcción requiere acceso a Internet: descarga las dependencias de Cargo
y de npm, la biblioteca qpdf, el código fuente de FFmpeg y la herramienta
que enlaza los módulos WASM.

\begin{longtable}{L{3.3cm} L{3.4cm} L{8.6cm}}
  \caption{Requisitos de software por etapa.}\label{tab:req-software} \\
  \toprule
  \textbf{Componente} & \textbf{Versión} & \textbf{Uso y fuente} \\
  \midrule
  \endfirsthead
  \toprule
  \textbf{Componente} & \textbf{Versión} & \textbf{Uso y fuente} \\
  \midrule
  \endhead
  \bottomrule
  \endfoot
    \multicolumn{3}{l}{\textbf{Construcción}} \\*
    \midrule
    Rust y Cargo
      & 1.88 o posterior
      & Compila el espacio de trabajo. Mínimo declarado en
        \texttt{rust-version} de \archivo{Cargo.toml}; se requiere
        \texttt{rustup} para añadir el destino
        \texttt{wasm32-unknown-unknown}. \\
    clang y Python 3
      & Sin versión fijada
      & Requeridos para generar los módulos Rust/WASM de la interfaz con
        \archivo{scripts/build-browser-crypto.sh} (clang y Python 3) y
        \archivo{scripts/build-evidence-verify.sh} (sin clang); véase
        \archivo{web/README.md}. \\
    Node.js y npm
      & Node.js 24
      & Instala dependencias y construye los portales de \archivo{web/}.
        Versión fijada en \archivo{.github/workflows/ci.yml} e indicada en
        \archivo{web/apps/qadra/README.md}. \\
    Motor de contenedores con Compose
      & Sin versión fijada
      & Opcional. \archivo{compose.yaml} define PostgreSQL 16 y Redis 7
        para desarrollo local; puede sustituirse por instancias propias de
        esas versiones. \\
    Biblioteca qpdf
      & 12.4.1
      & Validación de PDF. Se obtiene con
        \archivo{scripts/setup-document-formats.sh}, que verifica su SHA-256
        (\archivo{docs/document-format-operations.md}). \\
    FFmpeg y ffprobe
      & 9.0.2
      & Admisión de imagen, audio y video. Se compilan una sola vez con
        \archivo{scripts/install_media_decoder.py}, que requiere Python
        3.10 o posterior, GCC, make, NASM, pkg-config y los encabezados de
        zlib (\archivo{docs/media-decoder-setup.md}). \\
    OpenSSL y \texttt{unzip}
      & Sin versión fijada
      & Guiones de \archivo{pki/}, sellado de tiempo local y
        \archivo{scripts/demo.sh} (\archivo{README.md}). \\
    Herramientas de PostgreSQL y Redis
      & Sin versión fijada
      & El cliente \texttt{psql} se usa en la instalación local
        (\cref{cap:instalacion}). Con los ejecutables del servidor
        (\texttt{pg\_ctl}, \texttt{redis-server}, \texttt{redis-cli})
        \archivo{scripts/test-backends.sh} y \archivo{scripts/api-demo.sh}
        crean instancias desechables; el primero requiere además
        \texttt{setsid} (\archivo{README.md}). \\
    Git y \texttt{curl}
      & Sin versión fijada
      & Obtención del código y comprobación de la API en la instalación
        local y en \archivo{scripts/api-demo.sh}. \\
    TeX Live
      & Con LuaLaTeX, biber y makeglossaries
      & Compila el reporte (\archivo{latex/README.md}) y los manuales, que
        solo requieren LuaLaTeX y las fuentes Times New Roman y DejaVu Sans
        Mono (\archivo{manual-usuario/README.md}). \\
    \midrule
    \multicolumn{3}{l}{\textbf{Ejecución en servidor}} \\*
    \midrule
    Sistema operativo
      & Ubuntu 22.04 x86\_64
      & Instalación objetivo descrita en \archivo{docs/deployment.md}. El
        paquete debe compilarse sobre la misma distribución o una de ABI
        compatible. \\
    PostgreSQL
      & 16
      & Base de datos de la aplicación. \\
    Redis
      & 7
      & Desafíos, sesiones revocables y límites de uso; en el servidor se
        ejecuta con autenticación. \\
    nginx, OpenSSL y systemd
      & Sin versión fijada
      & nginx sirve la interfaz compilada y reenvía \texttt{/api} a la API;
        OpenSSL atiende a la CA y la TSA internas; las unidades de usuario
        de systemd supervisan los servicios bajo una cuenta sin privilegios
        de administrador. \\
    Python
      & 3.11 o posterior
      & Controlador de despliegue (\archivo{ops/}); no basta asumir la
        versión del intérprete del sistema. \\
    qpdf, FFmpeg y ffprobe
      & 12.4.1 y 9.0.2
      & No se instalan aparte: la biblioteca qpdf y los ejecutables
        verificados de FFmpeg y ffprobe acompañan al paquete de despliegue. \\
\end{longtable}

\section{Dependencias fijadas: qpdf y decodificador multimedia}%
\label{sec:req-fijadas}

Ambas dependencias son necesarias también en desarrollo para iniciar
\texttt{serve} o ejecutar las demostraciones con la API real: el servidor las
comprueba antes de abrir el puerto HTTP. Su preparación se describe en la
\cref{sec:inst-nativos}.

\importante{Un \texttt{ffmpeg} instalado desde los paquetes de la
distribución no sustituye a esta instalación: no garantiza el conjunto de
capacidades requerido. El prefijo predeterminado es
\archivo{/opt/tt-media}; si usa otro, prepárelo con el mismo
instalador e indíquelo con \texttt{TT\_FFMPEG\_PATH} y
\texttt{TT\_FFPROBE\_PATH}.}

\section{Recursos de hardware}\label{sec:req-hardware}

El repositorio no fija un dimensionamiento mínimo de procesador, memoria ni
almacenamiento. La única referencia disponible son los límites que la
provisión del servidor impone a cada servicio
(\archivo{docs/deployment.md}): 2~GiB de memoria para la API, 768~MiB para
PostgreSQL, 256~MiB para Redis y 128~MiB para nginx, con una cuota de dos
CPU por servicio. Son techos de consumo, no reservas ni un requisito mínimo
comprobado. El almacenamiento necesario depende del volumen de documentos
cifrados y de los respaldos que se conserven (\cref{cap:operacion}).

\section{Puertos en desarrollo}\label{sec:req-puertos}

La \cref{tab:req-puertos} lista los puertos predeterminados en un equipo de
desarrollo. Todos pueden cambiarse; los del servidor se asignan durante la
provisión y se describen en el \cref{cap:despliegue}.

\begin{table}[H]
  \centering
  \caption{Puertos predeterminados en desarrollo.}\label{tab:req-puertos}
  \begin{tabularx}{\textwidth}{L{1.6cm} L{4.2cm} Y}
    \toprule
    \textbf{Puerto} & \textbf{Servicio} & \textbf{Dónde se define} \\
    \midrule
    5432 & PostgreSQL & \archivo{compose.yaml}; \texttt{DATABASE\_URL} en
      \archivo{.env.example}. \\
    6379 & Redis & \archivo{compose.yaml}; \texttt{REDIS\_URL} en
      \archivo{.env.example}. \\
    3000 & API (\texttt{despacho-cli serve}) & Valor predeterminado de
      \texttt{-{}-bind}, \texttt{127.0.0.1:3000}. La interfaz reenvía
      \texttt{/api} a esa dirección salvo que se defina
      \texttt{API\_PROXY\_TARGET}
      (\archivo{web/apps/qadra/astro.config.mjs}). \\
    4321 & Interfaz Qadra (\texttt{npm run dev}) & Dirección indicada en
      \archivo{web/apps/qadra/README.md}; se cambia con
      \texttt{npm run dev -{}- -{}-port N} (\archivo{web/README.md}). \\
    18088 & Portal de la CA (\texttt{npm run dev:ca}) &
      \archivo{web/apps/ca/astro.config.mjs}. \\
    18089 & API de la CA & Destino predeterminado del proxy del portal de
      la CA; configurable con \texttt{CA\_API\_PROXY\_TARGET}
      (\archivo{docs/ca/operations.md}). \\
    \bottomrule
  \end{tabularx}
\end{table}

\section{Navegador del cliente}\label{sec:req-navegador}

Quien usa Qadra solo necesita un navegador web actualizado con JavaScript
activado; la firma y la validación de evidencia se ejecutan localmente en el
navegador mediante módulos Rust/WASM incluidos en la interfaz. Los
requisitos de cada persona usuaria (aplicación de autenticación para el
segundo factor, códigos de recuperación y certificado de la CA interna) se describen en el Manual de usuario,
capítulo «Antes de empezar».
